% รายงานการประชุมคณะกรรมการพัฒนาหลักสูตร ครั้งที่ 3/2566
% วาระ: ผลการสัมภาษณ์ผู้มีส่วนได้ส่วนเสียภายนอก 9 หน่วยงาน
% AUN-QA v4.0 — Criterion 1 (R1.4), Criterion 2 (R2.3, R2.6)
% หลักฐาน: AUNQA-1-4c
\documentclass{meetingminutes}

\meetingcommittee{คณะกรรมการพัฒนาหลักสูตร วิศวกรรมศาสตรมหาบัณฑิต
สาขาวิชาวิศวกรรมคอมพิวเตอร์และปัญญาประดิษฐ์ มหาวิทยาลัยราชภัฏอุตรดิตถ์}
\meetingnumber{๓}{๒๕๖๖}
\meetingdate{วันที่ ๒๐ พฤศจิกายน พ.ศ. ๒๕๖๖}
\meetingplace{ห้องประชุมคณะเทคโนโลยีอุตสาหกรรม}
\meetingstart{๑๓.๓๐ น.}
\meetingend{๑๖.๐๐ น.}

\begin{document}
\maketitle

\psection{ผู้มาประชุม}
\begin{participants}
  \pmember{1}{รองศาสตราจารย์ ดร.อิสระ อินจันทร์}{ประธานคณะกรรมการยกร่างหลักสูตร}{ประธาน}
  \pmember{2}{รองศาสตราจารย์ ดร.กันต์ อินทุวงศ์}{ที่ปรึกษาหลักสูตร}{รองประธาน}
  \pmember{3}{ผู้ช่วยศาสตราจารย์ ดร.อภิศักดิ์ พรหมฝาย}{อาจารย์ประจำหลักสูตร}{กรรมการ}
  \pmember{4}{ผู้ช่วยศาสตราจารย์ ดร.พิศิษฐ์ นาคใจ}{อาจารย์ประจำหลักสูตร}{กรรมการ}
  \pmember{5}{ผู้ช่วยศาสตราจารย์ ดร.พัชรี มณีรัตน์}{อาจารย์ประจำหลักสูตร}{กรรมการ}
  \pmember{6}{ผู้ช่วยศาสตราจารย์ ดร.พิทักษ์ คล้ายชม}{รองคณบดีฝ่ายวิชาการ คณะเทคโนโลยีอุตสาหกรรม}{กรรมการ}
  \pmember{7}{ผู้ช่วยศาสตราจารย์ ดร.โสกณ วิริยะรัตนกุล}{อาจารย์ประจำหลักสูตร}{กรรมการ}
  \pmember{8}{ผู้ช่วยศาสตราจารย์ ดร.ธัชชัย อยู่ยิ่ง}{อาจารย์ประจำหลักสูตร}{กรรมการ}
  \pmember{9}{ผู้ช่วยศาสตราจารย์ ดร.อภิศักดิ์ พรหมฝ่าย}{อาจารย์ประจำหลักสูตร}{กรรมการ}
  \pmember{7}{ผู้ช่วยศาสตราจารย์ ดร.กาญจนา ดาวเด่น}{อาจารย์ประจำหลักสูตร}{กรรมการ/เลขานุการ}
\end{participants}

\psection{ผู้ไม่มาประชุม}
\noindent ไม่มี\par

\startmeeting
\openingchair{รองศาสตราจารย์ ดร.อิสระ อินจันทร์}{กล่าวเปิดการประชุมและแจ้งให้ที่ประชุมทราบว่า
การประชุมครั้งนี้มีวาระสำคัญ คือ การรับฟังรายงานผลการสัมภาษณ์ผู้มีส่วนได้ส่วนเสียภายนอก
ซึ่ง ผู้ช่วยศาสตราจารย์ ดร.กาญจนา ดาวเด่น และ ผู้ช่วยศาสตราจารย์ ดร.พัชรี มณีรัตน์
ได้ดำเนินการสัมภาษณ์หน่วยงานต่าง ๆ ในจังหวัดอุตรดิตถ์และพื้นที่ใกล้เคียง จำนวน ๙ หน่วยงาน
ระหว่างเดือนตุลาคม–พฤศจิกายน ๒๕๖๖ เพื่อนำข้อมูลมาประกอบการพัฒนาหลักสูตร
จากนั้นดำเนินการประชุมตามระเบียบวาระ ดังนี้}

% ============================================================
\agenda{๑}{เรื่องแจ้งเพื่อทราบ}
% ============================================================

ประธานแจ้งให้ที่ประชุมทราบว่า คณะกรรมการมีมติในการประชุมครั้งที่ ๒/๒๕๖๖
ให้มอบหมาย ผู้ช่วยศาสตราจารย์ ดร.กาญจนา ดาวเด่น และ ผู้ช่วยศาสตราจารย์ ดร.พัชรี มณีรัตน์
ดำเนินการสัมภาษณ์ผู้มีส่วนได้ส่วนเสียภายนอกในพื้นที่จังหวัดอุตรดิตถ์
ซึ่งบัดนี้ดำเนินการแล้วเสร็จ และมีผลสรุปพร้อมนำเสนอต่อที่ประชุม

\noindent\textbf{มติที่ประชุม}\quad ที่ประชุมรับทราบ

% ============================================================
\agenda{๒}{รับรองรายงานการประชุมครั้งที่ ๒/๒๕๖๖}
% ============================================================

ที่ประชุมพิจารณารับรองรายงานการประชุมคณะกรรมการพัฒนาหลักสูตร ครั้งที่ ๒/๒๕๖๖
(วันที่ ๒๐ ตุลาคม ๒๕๖๖) ไม่มีกรรมการท่านใดขอแก้ไข

\resolution{ที่ประชุมมีมติรับรองรายงานการประชุมครั้งที่ ๒/๒๕๖๖ โดยไม่มีการแก้ไข}

% ============================================================
\agenda{๓}{รายงานผลการสัมภาษณ์ผู้มีส่วนได้ส่วนเสียภายนอก ๙ หน่วยงาน}
% ============================================================

ผู้ช่วยศาสตราจารย์ ดร.กาญจนา ดาวเด่น รายงานต่อที่ประชุมว่า คณะทำงานได้ดำเนินการ
สัมภาษณ์เชิงลึกกึ่งโครงสร้าง (Semi-structured Interview) กับผู้แทนองค์กรภายนอก
จำนวน ๙ หน่วยงาน โดยมีวัตถุประสงค์เพื่อ (๑) รวบรวมความต้องการของตลาดแรงงานและสังคม
ด้านวิศวกรรมคอมพิวเตอร์และปัญญาประดิษฐ์ในระดับภูมิภาค (๒) นำข้อมูลมาประกอบการกำหนด
ผลลัพธ์การเรียนรู้ของหลักสูตร และ (๓) ให้หลักสูตรตอบสนองบริบทท้องถิ่น
ใช้เวลาสัมภาษณ์ประมาณ ๔๕–๖๐ นาทีต่อหน่วยงาน ครอบคลุมประเด็น ๕ ด้าน ได้แก่
ทักษะและสมรรถนะที่ต้องการ ความคาดหวังต่อบัณฑิต ปัญหาที่พบในการปฏิบัติงาน
รูปแบบการศึกษาที่เหมาะสม และข้อเสนอแนะต่อหลักสูตร

\noindent รายชื่อหน่วยงานที่สัมภาษณ์ มีดังนี้

\vspace{6pt}
\noindent\small
\begin{tabular}{|>{\centering\arraybackslash}p{0.8cm}|p{5.5cm}|p{3.2cm}|p{4.0cm}|}
\hline
\bfseries ลำดับ & \bfseries หน่วยงาน & \bfseries ประเภท & \bfseries ผู้ให้สัมภาษณ์ \\
\hline
๑ & โรงพยาบาลอุตรดิตถ์ & สาธารณสุข/ภาครัฐ & ผู้แทนฝ่ายเทคโนโลยีสารสนเทศ \\
\hline
๒ & โรงพยาบาลพิชัย & สาธารณสุข/ภาครัฐ & ผู้แทนฝ่ายเทคโนโลยีสารสนเทศ \\
\hline
๓ & วิทยาลัยเทคนิคอุตรดิตถ์ & การศึกษา/อาชีวศึกษา & ผู้แทนฝ่ายวิชาการ (๒ ท่าน) \\
\hline
๔ & การไฟฟ้าส่วนภูมิภาคอุตรดิตถ์ & รัฐวิสาหกิจ/พลังงาน & ผู้แทนฝ่ายปฏิบัติการ \\
\hline
๕ & มหาวิทยาลัยเทคโนโลยีราชมงคลล้านนา พิษณุโลก & การศึกษา/อุดมศึกษา & ผู้แทนคณาจารย์สาขาวิศวกรรมคอมพิวเตอร์ \\
\hline
๖ & สหกรณ์ออมทรัพย์ครูอุตรดิตถ์ จำกัด & สถาบันการเงิน & ผู้แทนฝ่ายเทคโนโลยีสารสนเทศ \\
\hline
๗ & โรงเรียนเทศบาลท่าอิฐ & การศึกษา/โรงเรียน & ผู้แทนครูและฝ่ายวิชาการ \\
\hline
๘ & โรงเรียนเทศบาลวัดหนองผา & การศึกษา/โรงเรียน & ผู้แทนครูและฝ่ายวิชาการ \\
\hline
๙ & โรงเรียนเทศบาลวัดเกษมจิตตาราม & การศึกษา/โรงเรียน & ผู้แทนครูและฝ่ายวิชาการ \\
\hline
\end{tabular}
\normalsize

\vspace{8pt}
\noindent ผู้ช่วยศาสตราจารย์ ดร.กาญจนา ดาวเด่น สรุปผลการสัมภาษณ์รายหน่วยงาน ดังนี้

\subagenda{๓.๑}{โรงพยาบาลอุตรดิตถ์}

ต้องการบุคลากรที่สามารถวิเคราะห์ข้อมูลผู้ป่วยและนำ AI/Machine Learning มาช่วยระบุ
ผู้ป่วยความเสี่ยงสูง รวมถึงพัฒนาและดูแลระบบ Hospital Information System (HIS)
และ Electronic Medical Record (EMR) เน้นความสำคัญของความปลอดภัยข้อมูล
และกฎหมายคุ้มครองข้อมูลส่วนบุคคล (PDPA) ในบริบทข้อมูลทางการแพทย์

\noindent\textbf{PLO ที่สอดคล้อง:} PLO1 (การสร้างสรรค์ระบบ AI/ML), PLO2 (การประยุกต์ใช้ในองค์กร), PLO4 (จริยธรรมและ PDPA)

\subagenda{๓.๒}{โรงพยาบาลพิชัย}

ต้องการระบบดูแลผู้ป่วยระยะไกล (Telemedicine / Remote Patient Monitoring)
ที่สามารถใช้งานได้ในพื้นที่อินเทอร์เน็ตไม่เสถียร ต้องการผู้ที่สามารถวิเคราะห์ข้อมูล
เพื่อวางแผนสุขภาพชุมชน มีข้อเสนอแนะให้หลักสูตรมีโครงงานที่เชื่อมโยงกับบริบทท้องถิ่น

\noindent\textbf{PLO ที่สอดคล้อง:} PLO2 (การประยุกต์ใช้แก้ปัญหาชุมชน)

\subagenda{๓.๓}{วิทยาลัยเทคนิคอุตรดิตถ์}

ต้องการครู-อาจารย์ที่สามารถสอน AI และ Automation ในบริบทอาชีวศึกษา
ต้องการผู้เชี่ยวชาญช่วยพัฒนาระบบ Smart Factory / Industry 4.0 ในห้องปฏิบัติการ
และสนับสนุนให้ครูอาชีวศึกษาสามารถเข้าศึกษาต่อระดับบัณฑิตศึกษา

\noindent\textbf{PLO ที่สอดคล้อง:} PLO2 (การประยุกต์ใช้ในองค์กร), PLO3 (การสร้างองค์ความรู้), PLO5 (การถ่ายทอดและสื่อสาร)

\subagenda{๓.๔}{การไฟฟ้าส่วนภูมิภาคอุตรดิตถ์}

ต้องการบุคลากรที่วิเคราะห์ข้อมูล Smart Meter และระบบพลังงาน รวมถึงผู้เชี่ยวชาญ
ด้านความปลอดภัยของระบบไฟฟ้าและโครงสร้างพื้นฐาน (Cybersecurity / Critical Infrastructure)

\noindent\textbf{PLO ที่สอดคล้อง:} PLO1 (การสร้างสรรค์ระบบวิเคราะห์ข้อมูล), PLO4 (ความปลอดภัยและจริยธรรม)

\subagenda{๓.๕}{มหาวิทยาลัยเทคโนโลยีราชมงคลล้านนา พิษณุโลก}

ต้องการความร่วมมือด้านวิจัย AI / Computer Engineering ระหว่างสถาบัน
และการแลกเปลี่ยนนักศึกษาระดับบัณฑิตศึกษา เพื่อสร้างเครือข่ายวิชาการในภูมิภาค

\noindent\textbf{PLO ที่สอดคล้อง:} PLO3 (การสร้างองค์ความรู้), PLO5 (การทำงานร่วมกัน/เครือข่าย)

\subagenda{๓.๖}{สหกรณ์ออมทรัพย์ครูอุตรดิตถ์ จำกัด}

ต้องการนำ AI มาวิเคราะห์ความเสี่ยงสินเชื่อ (Credit Risk Analysis) และตรวจจับธุรกรรมผิดปกติ
ต้องการผู้เชี่ยวชาญด้าน Cybersecurity และ Data Protection เพื่อปกป้องข้อมูลทางการเงิน
เน้นความเข้าใจกฎหมายคุ้มครองข้อมูลส่วนบุคคล (PDPA)

\noindent\textbf{PLO ที่สอดคล้อง:} PLO1 (ระบบ AI วิเคราะห์ความเสี่ยง), PLO2 (การประยุกต์ใช้ในองค์กร), PLO4 (PDPA และ Cybersecurity)

\subagenda{๓.๗–๓.๙}{โรงเรียนเทศบาล ๓ แห่ง (ท่าอิฐ, วัดหนองผา, วัดเกษมจิตตาราม)}

ต้องการครูที่ใช้ AI สร้างสื่อการสอนและพัฒนานวัตกรรมการเรียนรู้ในยุคดิจิทัล
ต้องการระบบ E-Portfolio เพื่อสนับสนุนการเลื่อนวิทยฐานะของครู
และต้องการบุคลากรที่สามารถถ่ายทอดความรู้ด้าน AI ให้กับครูในโรงเรียนได้อย่างเป็นรูปธรรม
มีข้อเสนอแนะให้มีกิจกรรม Experiential Learning ที่นักศึกษาถ่ายทอดความรู้สู่ชุมชน

\noindent\textbf{PLO ที่สอดคล้อง:} PLO1 (การสร้างสรรค์สื่อ/แพลตฟอร์ม), PLO2 (การประยุกต์ใช้ในโรงเรียน), PLO5 (การถ่ายทอดและสื่อสาร)

\resolution{ที่ประชุมมีมติรับทราบรายงานผลการสัมภาษณ์ผู้มีส่วนได้ส่วนเสียภายนอก จำนวน ๙ หน่วยงาน
และขอบคุณ ผู้ช่วยศาสตราจารย์ ดร.กาญจนา ดาวเด่น และ ผู้ช่วยศาสตราจารย์ ดร.พัชรี มณีรัตน์
ที่ดำเนินการสัมภาษณ์อย่างครบถ้วน}

% ============================================================
\agenda{๔}{พิจารณาความสอดคล้องของ PLO กับความต้องการผู้มีส่วนได้ส่วนเสีย
และแนวทางปรับปรุงหลักสูตร}
% ============================================================

ผู้ช่วยศาสตราจารย์ ดร.กาญจนา ดาวเด่น นำเสนอตาราง Needs-to-PLO Mapping
ซึ่งแสดงให้เห็นความสอดคล้องระหว่างความต้องการของผู้มีส่วนได้ส่วนเสียกับ PLO ทั้ง ๕ ข้อ

\vspace{6pt}
\noindent\small
\begin{tabular}{|p{5.0cm}|p{2.8cm}|>{\centering\arraybackslash}p{1.0cm}|>{\centering\arraybackslash}p{1.0cm}|>{\centering\arraybackslash}p{1.0cm}|>{\centering\arraybackslash}p{1.0cm}|>{\centering\arraybackslash}p{1.0cm}|}
\hline
\bfseries ความต้องการหลัก & \bfseries แหล่งที่มา & \bfseries PLO1 & \bfseries PLO2 & \bfseries PLO3 & \bfseries PLO4 & \bfseries PLO5 \\
\hline
วิเคราะห์ข้อมูล/พัฒนาระบบ AI-ML สำหรับงานเฉพาะทาง
  & รพ.อุตรดิตถ์, รพ.พิชัย, กฟภ., สหกรณ์ฯ & $\bullet$ & & $\bullet$ & & \\
\hline
ประยุกต์ใช้ความรู้แก้ปัญหาองค์กรและชุมชนท้องถิ่น
  & ทุกหน่วยงาน & & $\bullet$ & & & \\
\hline
AI/Automation สำหรับอาชีวศึกษาและ Smart Factory
  & ว.เทคนิคฯ & & $\bullet$ & & & \\
\hline
Cybersecurity และจริยธรรมข้อมูล (PDPA)
  & กฟภ., สหกรณ์ฯ, รพ. & & & & $\bullet$ & \\
\hline
วิจัยและสร้างองค์ความรู้ใหม่ (วิทยานิพนธ์/สารนิพนธ์)
  & มทร.ล้านนาฯ, ว.เทคนิคฯ & & & $\bullet$ & & \\
\hline
AI สำหรับการสร้างสื่อการสอนและ E-Portfolio
  & รร.เทศบาล ๓ แห่ง & $\bullet$ & $\bullet$ & & & $\bullet$ \\
\hline
ทักษะถ่ายทอดความรู้ สื่อสาร ทำงานเป็นทีม
  & รร.เทศบาล, มทร.ล้านนาฯ & & & & & $\bullet$ \\
\hline
ความร่วมมือวิจัยระหว่างสถาบัน
  & มทร.ล้านนา พิษณุโลก & & & $\bullet$ & & $\bullet$ \\
\hline
\end{tabular}
\normalsize

\vspace{8pt}
\noindent ที่ประชุมร่วมกันอภิปราย มีประเด็นสำคัญดังนี้

\begin{thailist}
  \item ผู้ช่วยศาสตราจารย์ ดร.พิศิษฐ์ นาคใจ เสนอว่า ควรเพิ่มเนื้อหาด้าน Cybersecurity
        และ PDPA ให้ชัดเจนใน CLO ของรายวิชาที่เกี่ยวข้อง เพื่อให้ตรงกับความต้องการของ
        หน่วยงานด้านสาธารณสุข พลังงาน และสถาบันการเงิน ที่ประชุมเห็นชอบ
  \item รองศาสตราจารย์ ดร.กันต์ อินทุวงศ์ เสนอว่า ควรออกแบบสารนิพนธ์และวิทยานิพนธ์
        ให้เชื่อมโยงกับปัญหาจริงในองค์กรและชุมชนท้องถิ่น เพื่อตอบสนอง PLO2
        และสอดคล้องกับปรัชญามหาวิทยาลัยราชภัฏ ที่ประชุมเห็นชอบ
  \item ผู้ช่วยศาสตราจารย์ ดร.พัชรี มณีรัตน์ เสนอว่า ควรบรรจุกิจกรรม Experiential Learning
        ที่ให้นักศึกษาถ่ายทอดความรู้ด้าน AI สู่โรงเรียนและชุมชน
        ตามที่โรงเรียนเทศบาล ๓ แห่งระบุ ซึ่งสนับสนุน PLO5 ที่ประชุมเห็นชอบ
\end{thailist}

\resolution{ที่ประชุมมีมติเห็นชอบแนวทางปรับปรุงหลักสูตร ดังนี้
(๑) PLO ทั้ง ๕ ข้อมีความสอดคล้องกับความต้องการของผู้มีส่วนได้ส่วนเสียในระดับดี
ให้คงไว้ตามที่ยกร่าง
(๒) มอบหมายให้ ผู้ช่วยศาสตราจารย์ ดร.พิศิษฐ์ นาคใจ นำข้อมูล Needs-to-PLO Mapping
ไปประกอบการยกร่าง CLO ของรายวิชา โดยเพิ่มเนื้อหา Cybersecurity และ PDPA ให้ชัดเจน
ภายในเดือนธันวาคม ๒๕๖๖
(๓) มอบหมายให้ ผู้ช่วยศาสตราจารย์ ดร.อภิศักดิ์ พรหมฝาย ออกแบบรูปแบบสารนิพนธ์/วิทยานิพนธ์
ที่เชื่อมโยงกับปัญหาจริงในองค์กรท้องถิ่น เพื่อนำเสนอในการประชุมครั้งถัดไป
(๔) มอบหมายให้ ผู้ช่วยศาสตราจารย์ ดร.กาญจนา ดาวเด่น จัดทำข้อเสนอกิจกรรม
Experiential Learning ในรายวิชาสัมมนาและหัวข้อพิเศษ เพื่อให้นักศึกษาถ่ายทอดความรู้สู่ชุมชน}

% ============================================================
\agenda{๕}{เรื่องอื่น ๆ}
% ============================================================

ผู้ช่วยศาสตราจารย์ ดร.พิทักษ์ คล้ายชม แจ้งให้ที่ประชุมทราบว่า คณะเทคโนโลยีอุตสาหกรรม
พร้อมสนับสนุนทรัพยากรสำหรับการจัดการเรียนการสอน และขอให้คณะกรรมการพัฒนาหลักสูตร
เร่งดำเนินการยกร่างเอกสาร มคอ.๒ เพื่อเสนอผ่านระดับคณะภายในปีการศึกษา ๒๕๖๖

รองศาสตราจารย์ ดร.อิสระ อินจันทร์ กำหนดการประชุมครั้งถัดไป คือ ครั้งที่ ๔/๒๕๖๖
ในเดือนธันวาคม ๒๕๖๖ เพื่อพิจารณา PLO ฉบับสมบูรณ์และโครงสร้างรายวิชา

\noindent\textbf{มติที่ประชุม}\quad ที่ประชุมรับทราบ

\adjournmeeting

\begin{sigblock}
\typist{ผู้ช่วยศาสตราจารย์ ดร.กาญจนา ดาวเด่น}
\reviewer{ผู้ช่วยศาสตราจารย์ ดร.พิศิษฐ์ นาคใจ}{กรรมการ}
\chairsig{รองศาสตราจารย์ ดร.อิสระ อินจันทร์}{ประธานที่ประชุม}
\end{sigblock}

\end{document}
